% !TEX spellcheck = de
\DocumentMetadata{lang=de}
\documentclass{ltx-talk}

\makeatletter
\providecommand*{\input@path}{}
\g@addto@macro\input@path{{./include/}{./beamertheme-sam/}}
\makeatother

\definecolor{theme}{HTML}{870030}

\usepackage[
  themecolor=theme,
]{talkthemesam}

\usepackage[ngerman]{babel}
\usepackage{csquotes}
\usepackage{fontawesome7}
\usepackage{tikz}
\usetikzlibrary{shapes.geometric}
\usepackage[style=authoryear]{biblatex}
\addbibresource{ltx-talk.bib}
\DeclareLabeldate{%
  \field{date}
  \field{year}
  \field{eventdate}
  \field{origdate}
  \field{urldate}
}
\NewDocumentEnvironment{notes}{b}{}{}

\title{Von Beamer zu ltx-talk}
\subtitle{Ein Nutzerbericht}
\titlegraphic{Schloss_Biebrich,_Nordseite,_Wiesbaden-Biebrich,_150904,_ako}
\titlecredit{
  Titelbild: \href{https://commons.wikimedia.org/wiki/File:Schloss_Biebrich,_Nordseite,_Wiesbaden-Biebrich,_150904,_ako.jpg}{Ansgar Koreng (CC BY 3.0 via wikimedia commons)}%
}
\author{samcarter}
\institute{Dante Herbsttagung Wiesbaden}
\date{25.\,-- 26.\ September 2026}

\begin{document}

\maketitle
\begin{notes}
Danke für die Einführung. Wie schon erwähnt, werde ich heute über den Wechsel von beamer zu ltx-talk sprechen aus der Perspektive eines Nutzers sprechen.
\end{notes}

\begin{frame}
\frametitle{Was ist ltx-talk}

\begin{notes}
content...
\end{notes}
\end{frame}

\begin{frame}
  \thispagestyle{plain}
  \begin{tikzpicture}[
    remember picture,
    overlay,
    bubble/.style={
      fill=white,
      fill opacity=0.7,
      text opacity=1,
      inner sep=0.3cm,
      font=\Large,
      shape=ellipse
    }
  ]
    \node at (current page.center) {\includegraphics[width=\paperwidth]{David_Garrick_Between_Tragedy_and_Comedy}};
    \node[bubble] at (3,-3) {ltx-talk};
    \node[bubble] at (7,0) {me};
    \node[bubble] at (12.5,-2) {beamer};
    \node[anchor=south,white] at (current page.south) {Joshua Reynolds (1761) \emph{\enquote{David Garrick zwischen Tragödie und Komödie}}};
  \end{tikzpicture}
  \begin{notes}
    so, was hat mich nun dazu bewegt, von beamer, von dem ich schon seit Jahren ein großer Fan bin, zu ltx-talk zu wechseln? 
  \end{notes}
\end{frame}

%\begin{frame}<-2>
%  \frametitle{Zisterzienser Ziffern}
%  \begin{minipage}{.48\textwidth}
%    \begin{onlyenv}<1>%
%      \begin{itemize}
%        \item mittelalterliches Zahlensystem
%        \item Zisterzienser Orden, frühes 13. bis 15. Jhd
%        \item andere Verwendungen bis ins frühe 18. Jhd
%        \begin{itemize}
%        \item Weinhandel um Brügge
%        \item Freimaurer in Paris
%        \item \ldots
%        \end{itemize}
%      \end{itemize}
%    \end{onlyenv}%
%    \begin{onlyenv}<2->%
%      \cistercianset{height=4cm,stroke=2mm}
%      \begin{tikzpicture}
%        \node (cis) {%
%          \only<-2>{\cistercian{1337}}%
%          \only<3->{\cistercian{1989}}%
%        };
%
%        \node[samorange,font=\bfseries,anchor=south west] at (1.5,1) {Einer};
%        \draw<2>[samorange,line width=2mm] (0.1,1.9) -- ++(1.2333,0) -- ++(0,-1.3333);
%        \draw<3->[samorange,line width=2mm] (0.1,1.9) -| ++(1.2333,-1.3333) -- ++(-1.2333,0);
%
%        \node[samred,font=\bfseries,anchor=south east] at (-1.5,1) {Zehner};
%        \draw<2>[samred,line width=2mm] (-0.1,1.9) -- ++(-1.3333,-1.3333);
%        \draw<4->[samred,line width=2mm] (-0.1,2) ++ (-1.2333,0) |- ++(1.2333,-1.4333);
%
%        \node[samlblue,font=\bfseries,anchor=north west] at (1.5,-1) {Hunderter};
%        \draw<2>[samlblue,line width=2mm] (0.1,-1.9) -- ++(1.3333,1.3333);
%        \draw<5->[samlblue,line width=2mm] (0.1,-1.9) -| ++(1.2333,1.3333) -- ++(-1.2333,0);
%
%        \node[samlgreen,font=\bfseries,anchor=north east] at (-0.9,-1) {Tausender};
%        \draw<2,6->[samlgreen,line width=2mm] (-0.1,-1.9) -- ++(-1.3333,0);
%
%        \draw<2>[black,line width=2mm] (0,-2) -- (0,2);
%
%        \node[anchor=north,font=\bfseries\Huge] at ([yshift=-1cm]cis.south) {-\kern0.15em-\kern0.15em-\kern0.15em-};
%        \node<2>[anchor=north,font=\bfseries\Huge] at ([yshift=-0.3cm]cis.south) {%
%          \textcolor{samlgreen}{1}%
%          \textcolor{samlblue}{3}%
%          \textcolor{samred}{3}%
%          \textcolor{samorange}{7}%
%        };
%        \node<2>[anchor=north,overlay] at ([yshift=-1.5cm]cis.south) {(MCCCXXXVII)};
%        \node<3->[anchor=north,font=\bfseries\Huge] at ([yshift=-0.3cm]cis.south) {%
%          \uncover<7->{\textcolor{samlgreen}{1}}%
%          \uncover<6->{\textcolor{samlblue}{9}}%
%          \uncover<5->{\textcolor{samred}{8}}%
%          \uncover<4->{\textcolor{samorange}{9}}%
%        };
%      \end{tikzpicture}%
%    \end{onlyenv}%
%  \end{minipage}%
%  \hfill  
%  \begin{minipage}{.5\textwidth}
%    \cistercianset{height=1cm,stroke=1mm,baseline=0.3cm}%
%    \setcounter{mycell}{0}%
%    \begin{tblr}{
%      width=.9\textwidth,
%      colspec={@{}XXXX[2]@{}},
%      cells={
%        halign=r,
%        cmd={%
%          \addtocounter{mycell}{1}%
%          \themycell:%
%          \cistercian{\themycell}%
%        }
%      },
%      cell{1-3}{4}={
%        cmd={
%          \fpeval{10^\therownum}:%
%          \cistercian{\fpeval{10^\therownum}}%
%        }
%      },
%      rows={rowsep+=2mm},
%    }
%      & & &\\
%      & & &\\
%      & & &\\
%    \end{tblr}
%    
%    \bigskip
%    
%    (nach einem Astrolabium aus dem frühen 16. Jhd)
%  \end{minipage}%
%\end{frame}
%
%\begin{frame}
%  \frametitle{\LaTeX\ Pakete}
%  \label{packages}
%  \begin{columns}[vertical-alignment=top]
%    \begin{column}[t]{.5\textwidth}
%      {\Huge\packagex\par}
%      \bigskip
%      \colorlet{structure}{samred}
%      \begin{itemize}
%        \item Autor: Jonathan Spratte
%        \item sehr schnell
%        \item viele Optionen
%        \item Querformat und alt.\ Formen
%      \end{itemize}
%    \end{column}
%    \begin{column}{.45\textwidth}
%      {\Huge\packagec\par}
%      \bigskip
%      \colorlet{structure}{samlblue}
%      \begin{itemize}
%        \item Autor: samcarter
%        \item Ti\emph{k}Z-basiert
%        \item Linien-Übergänge
%        \item Anpassung an Schriftart
%      \end{itemize}
%    \end{column}
%  \end{columns}
%\end{frame}
%
%\IfFileExists{Xisterzienser.pdf}{\includepdf[pages=1-]{Xisterzienser}}{}
%
%\begingroup
%\colorlet{structure}{samlblue}
%\begin{frame*}
%  \frametitle{\packagec*: Einfaches Beispiel}
%  \begin{tcblisting}{}
%    \documentclass{article}
%
%    \usepackage{cistercian}
%    \pagenumbering{cistercian}
%
%    \begin{document}
%
%    \cistercian{3} Haselnüsse für Aschenbrödel
%
%    \end{document}
%  \end{tcblisting}
%\end{frame*}
%
%\begin{frame*}
%  \frametitle{\packagec*: Anpassungsfähig}
%  \begin{tcblisting}{}
%    \color{red} Jim Knopf und die wilde \cistercian{13}
%  \end{tcblisting}
%  \begin{tcblisting}{}
%    \Huge Die \cistercian{7} Zwerge
%  \end{tcblisting}
%  \begin{tcblisting}{}
%    \Huge\fontfamily{ugq}\selectfont
%    \cistercian{1001} Nacht
%  \end{tcblisting}
%\end{frame*}
%
%\begin{frame*}
%  \frametitle{\packagec*: Optionen}
%  \begin{tcblisting}{}
%    \cistercianset{stroke=1mm}
%
%    Ali Baba und die \cistercian[height=1cm]{40} Räuber
%  \end{tcblisting}
%\end{frame*}
%
%\begin{frame*}
%  \frametitle{\packagec*: Ti\emph{k}Z-basiert}
%  \begin{tcblisting}{}
%    Goldlöckchen und die \cistercian[draw=red]{3} Bären
%  \end{tcblisting}
%  \begin{tcblisting}{before lower app=\hskip-1cm}
%    \cistercianset{scale=8}
%    \cistercian[line join=round]{3}
%    \cistercian[line join=bevel]{3}
%    \cistercian[line join=miter]{3}
%  \end{tcblisting}
%\end{frame*}
%\endgroup
%
%\begin{frame*}
%  \frametitle{\protect\tikz[scale=0.6,baseline=0.5cm]{\panda}~~Panda Pakete}
%  \begin{tcblisting}{
%    before lower app=\setlength{\parskip}{\bigskipamount},
%    listing and comment,
%    listing outside comment,
%    comment={%
%      \fontfamily{cmr}\selectfont
%      Computer Modern Roman:
%      \rule{ \CleverPanda{} }{1.6ex}
%  
%      \fontfamily{ugq}\selectfont
%      URW Grotesk:
%      \rule{ \CleverPanda{} }{1.6ex}
%    }
%  }
%    \usepackage{CleverPanda}
%
%    \rule{ \CleverPanda{} }{1.6ex}
%  \end{tcblisting}
%
%  \bigskip
%
%  \begin{tcblisting}{
%    before lower app=\setlength{\parskip}{\bigskipamount},
%    listing and comment,
%    listing outside comment,
%    comment={%
%      \fontfamily{cmr}\selectfont
%      Computer Modern Roman:
%      \BusyPanda{}
%  
%      \fontfamily{ugq}\selectfont
%      URW Grotesk:
%      \BusyPanda{}
%    }
%  }
%    \usepackage{BusyPanda}
%
%    \BusyPanda{}
%  \end{tcblisting}
%\end{frame*}
%
\begin{frame}
  \frametitle{Abenteuerlustig?}
  \begin{itemize}
    \item \fullcite{ltx-talk}
    \item \fullcite{Wright:TB47-1:64}
    \item Tim Prescott (2026). \enquote{Creating accessible presentations with ltx-talk}. TUG 2026 (Calgary). \textsc{recording:} \url{https://www.youtube.com/watch?v=pltqpUBeG24}
    \linebreak\textsc{preprint:} \url{https://tug.org/tug2026/preprints/prescott-ltx-talk.pdf}
  \end{itemize}
\end{frame}


\end{document}